\documentclass[12pt,a4paper]{article}

% ------------------------------------------------
% Pacotes
% ------------------------------------------------
\usepackage[utf8]{inputenc}
\usepackage[T1]{fontenc}
\usepackage[brazil]{babel}
\usepackage{graphicx}
\usepackage[table]{xcolor}
\usepackage{array}
\usepackage{tabularx}
\usepackage{geometry}
\usepackage{fancyhdr}
\usepackage{mathptmx} % Times New Roman
\usepackage{enumitem}\usepackage{amsfonts}
\usepackage{amsmath} % Recomendado para o comando \text{}
\usepackage{amssymb}
\usepackage{tikz}
\usetikzlibrary{positioning, shapes.geometric}

% TikZ Styles for ER Diagrams
\tikzset{
    block/.style={rectangle, draw, thick, minimum width=2cm, minimum height=1cm, fill=white},
    diamond_node/.style={diamond, draw, thick, minimum width=2cm, minimum height=1cm, fill=white, aspect=2},
    attribute/.style={ellipse, draw, thin, minimum width=1.5cm, minimum height=0.6cm, fill=white, font=\scriptsize},
    line/.style={draw, thick},
}
% ------------------------------------------------ 
% Configuração de página
% ------------------------------------------------
\geometry{
    a4paper,
    left=2cm,
    right=2cm,
    top=2cm,
    bottom=2cm,
    headheight=15pt 
}% Caixa para resposta discursiva
\newcommand{\resposta}[1]{
\noindent
\framebox[\textwidth]{
\begin{minipage}{0.98\textwidth}
\vspace{#1}
\end{minipage}
}
}

\pagestyle{fancy}
\fancyhf{}
\renewcommand{\headrulewidth}{0pt}

\setlength{\parindent}{0pt}

% ------------------------------------------------
% Comandos auxiliares
% ------------------------------------------------
\newcolumntype{L}[1]{>{\raggedright\arraybackslash}p{#1}}
\newcolumntype{C}[1]{>{\centering\arraybackslash}p{#1}}

% ------------------------------------------------
% Documento
% ------------------------------------------------
\begin{document}

% ------------------------------------------------
% Cabeçalho
% ------------------------------------------------
\begin{tabularx}{\textwidth}{@{} L{4cm} X @{}}
    \raisebox{-0.5cm}{\includegraphics[width=4cm]{Universidade_de_Rio_Verde.jpg}} &
    \hfill
    \begin{tabular}{|>{\centering\arraybackslash}m{10cm}|>{\raggedright\arraybackslash}p{1.8cm}|}
        \hline
        \cellcolor{black}\color{white}\rule{0pt}{1.2cm}\textbf{AVALIAÇÃO 2 - ADAPTADA} 
        & 
        \scriptsize Nota \tabularnewline
        \hline
    \end{tabular} 
\end{tabularx}
\vspace{0.3cm}

% ------------------------------------------------  
% Tabela principal
% ------------------------------------------------
\begin{tabularx}{\textwidth}{|X|L{2cm}|}
\hline
\textbf{Curso} & \textbf{Campus} \\
\textit{Engenharia de Software} & \textit{Rio Verde} \\
\hline
\multicolumn{2}{|l|}{\textbf{Disciplina}} \\  
\multicolumn{2}{|l|}{ \textit{Banco de dados}} \\
\hline
\textbf{Nome do(a) acadêmico(a)} & \textbf{Valor} \\

MARCIO HENRIQUE LENZ SILVA& 7,0 \\
\hline
\textbf{Professor} & \textbf{Turma} \\
\textit{ Sergio Souza Novak} & \textit{B} \\
\hline
\textbf{Data da Avaliação} & \textbf{Código} \\
\textit{12/05/2026} & \textit{ESW417} \\
\hline
\end{tabularx}
 
% ------------------------------------------------
% Barra de Atenção
% ------------------------------------------------
\vspace{0.2cm}

\begin{tabularx}{\textwidth}{|X|}
\hline
\cellcolor{black}\color{white}
\centering \textbf{ATENÇÃO:} Avaliação adaptada para alunos com deficiência auditiva. \tabularnewline
\hline
\end{tabularx}

\vspace{0.3cm}

\textbf{ORIENTAÇÕES PARA A RESOLUÇÃO}  

\begin{itemize}
    \item A avaliação é individual, sem consulta;
    \item As questões teóricas são de múltipla escolha;
    \item A parte prática pode ser enviada em PDF por e-mail.
\end{itemize}

\vspace{0.5cm}

\begin{center}
    \rule{\textwidth}{0.4pt}
    \textbf{QUESTÕES TEÓRICAS}
\end{center}

\vspace{0.2cm}

\textbf{INSTRUÇÕES:} Esta seção contém \textbf{5 exercícios}. Recomenda-se reservar \textbf{40 minutos} para sua resolução. Nesta etapa teórica, \textbf{não é permitido} o uso de aparelhos celulares nem o acesso a qualquer ferramenta de Inteligência Artificial. \textbf{Todas as questões são de assinalar.}

\vspace{0.3cm}

\begin{enumerate}
    \item \textbf{(Ano: 2026 Banca: Instituto Legalle Órgão: BADESUL - RS Prova: Analista de Sistemas)} \\
    Os bancos de dados relacionais se baseiam em princípios que garantem integridade, consistência e organização dos dados.
    
    Analise as afirmativas:
    
    \begin{enumerate}[label=\Roman*.]
        \item Uma chave primária pode ser composta por mais de um atributo, sendo denominada chave primária composta.
        \item Uma chave estrangeira em uma tabela deve referenciar obrigatoriamente a chave primária de uma tabela diferente, nunca podendo referenciar a própria tabela.
        \item A normalização é um processo que visa reduzir a redundância de dados e eliminar anomalias de inserção, atualização e exclusão.
        \item Na Primeira Forma Normal (1FN), cada atributo de uma tupla deve conter apenas valores atômicos, ou seja, indivisíveis.
    \end{enumerate}
    
    Estão corretas:
    
    \begin{enumerate}[label=\Alph*)]
        \item I e II, apenas.
        \item II e III, apenas.
        \item I, II e IV, apenas.
        \item I, III e IV, apenas.
        \item I, II, III e IV.
    \end{enumerate}

    \item A tabela abaixo viola a Primeira Forma Normal (1FN). Assinale a alternativa que identifica corretamente o \textbf{motivo da violação} e a \textbf{normalização correta}:
    
    \begin{center}
        \begin{tabular}{|c|c|L{8cm}|}
            \hline
            \cellcolor{gray!20}\textbf{ID} & \cellcolor{gray!20}\textbf{Nome} & \cellcolor{gray!20}\textbf{Contatos} \\ \hline
            1 & Carlos & car@email.com, (11) 9111-1111 \\ \hline
            2 & Ana & ana@email.com \\ \hline
            3 & Bia & bia@work.com, 98460-1364 \\ \hline
        \end{tabular}
    \end{center}

    \begin{enumerate}[label=\Alph*)]
        \item Motivo: Falta de chave primária. Normalização: Definir o ID como PK.
        \item Motivo: Atributos multivalorados (campo Contatos possui email e telefone). Normalização: Criar um registro para cada contato.
        \item Motivo: Dependência parcial da chave. Normalização: Criar uma tabela para o Nome.
        \item Motivo: Dependência transitiva. Normalização: Remover o campo Nome.
        \item Motivo: Redundância de IDs. Normalização: Alterar o ID da Bia para 4.
    \end{enumerate}


    \item A tabela \textbf{VENDAS} registra vendas de uma farmácia. Considere que a Chave Primária é o atributo \textbf{ID\_Venda}.

    \begin{center}
        \small
        \begin{tabular}{|c|c|c|l|c|}
            \hline
            \cellcolor{gray!20}\textbf{ID\_Venda} & \cellcolor{gray!20}\textbf{Data} & \cellcolor{gray!20}\textbf{ID\_Prod} & \cellcolor{gray!20}\textbf{Nome\_Prod} & \cellcolor{gray!20}\textbf{Preco\_Unit} \\ \hline
            V001 & 01/03 & P10 & Dipirona 500mg & R\$ 4,50 \\ \hline
            V002 & 01/03 & P20 & Amoxicilina & R\$ 18,00 \\ \hline
            V003 & 02/03 & P10 & Dipirona 500mg & R\$ 4,50 \\ \hline
            V004 & 02/03 & P30 & Ibuprofeno & R\$ 12,00 \\ \hline
            V005 & 03/03 & P20 & Amoxicilina & R\$ 18,00 \\ \hline
        \end{tabular}
    \end{center}

    Julgue as afirmativas como Verdadeiras (V) ou Falsas (F):\\
         ( \quad ) A tabela está na Primeira Forma Normal (1FN).\\
         ( \quad ) A tabela está na Segunda Forma Normal (2NF).\\
         ( \quad ) A tabela está na Terceira Forma Normal (3NF).\\
         ( \quad ) A tabela está na Forma Normal de Boyce-Codd (FNBC).\\


    \item Em uma empresa, cada representante atende a várias regiões e vende várias linhas de produtos. Sabe-se que as regiões atendidas e as linhas de produtos vendidas são \textbf{independentes} entre si. Todos os atributos compõem a Chave Primária.
    
    \begin{center}
        \small
        \begin{tabular}{|l|l|l|}
            \hline
            \cellcolor{gray!20}\textbf{Representante (PK)} & \cellcolor{gray!20}\textbf{Região (PK)} & \cellcolor{gray!20}\textbf{Linha\_Produto (PK)} \\ \hline
            Carlos Silva & Norte & Máquinas \\ \hline
            Carlos Silva & Norte & EPIs \\ \hline
            Carlos Silva & Sul & Ferramentas \\ \hline
            Carlos Silva & Sul & Máquinas \\ \hline
            Carlos Silva & Sul & EPIs \\ \hline
            Carlos Silva & Leste & Ferramentas \\ \hline
            Carlos Silva & Leste & Máquinas \\ \hline
            Carlos Silva & Leste & EPIs \\ \hline
        \end{tabular}
    \end{center}

    Assinale a alternativa que identifica a \textbf{Forma Normal violada} e a \textbf{decomposição correta} para atingir a \textbf{Quarta Forma Normal (4FN)}:

    \begin{enumerate}[label=\Alph*)]
        \item Viola a 2FN. Decomposição: (Representante, Região) e (Região, Linha).
        \item Viola a 3FN. Decomposição: (Representante, Linha) e (Região, Linha).
        \item Viola a 4FN (Dependência Multivalorada). Decomposição: (Representante, Região) e (Representante, Linha\_Produto).
        \item Viola a FNBC. Decomposição: Tabela Única.
        \item Viola a 1FN. Decomposição: (Representante, Nome).
    \end{enumerate}

    \item Os diagramas de Venn abaixo representam visualmente as operações de \textbf{JOIN} entre duas tabelas, \textbf{A} (esquerda) e \textbf{B} (direita). Assinale a alternativa que identifica corretamente os tipos de JOIN representados nas Figuras I e II:

    \begin{center}
        \begin{tikzpicture}[thick]
            % Figura I: INNER JOIN
            \begin{scope}
                \clip (0,0) circle (1cm);
                \fill[gray!40] (1.2,0) circle (1cm);
            \end{scope}
            \draw (0,0) circle (1cm) node[xshift=-0.5cm] {A};
            \draw (1.2,0) circle (1cm) node[xshift=0.5cm] {B};
            \node at (0.6, -1.5) {\textbf{Figura I}};
        \end{tikzpicture}
        \hspace{2.5cm}
        \begin{tikzpicture}[thick]
            % Figura II: LEFT JOIN
            \fill[gray!40] (0,0) circle (1cm);
            \draw (0,0) circle (1cm) node[xshift=-0.5cm] {A};
            \draw (1.2,0) circle (1cm) node[xshift=0.5cm] {B};
            \node at (0.6, -1.5) {\textbf{Figura II}};
        \end{tikzpicture}
    \end{center}

    \begin{enumerate}[label=\Alph*)]
        \item Figura I: LEFT JOIN; Figura II: INNER JOIN.
        \item Figura I: INNER JOIN; Figura II: LEFT JOIN.
        \item Figura I: RIGHT JOIN; Figura II: FULL JOIN.
        \item Figura I: INNER JOIN; Figura II: RIGHT JOIN.
        \item Figura I: FULL JOIN; Figura II: LEFT JOIN.
    \end{enumerate}

\end{enumerate}
\vspace{0.5cm}
\newpage
\begin{center}
    \rule{\textwidth}{0.4pt}
    \textbf{QUESTÕES PRÁTICAS}
\end{center}

\vspace{0.2cm}

\textbf{INSTRUÇÕES:} Esta seção contém \textbf{5 exercícios}. Para acadêmicos com deficiência auditiva, é facultado o envio da resolução em formato \textbf{PDF} para o e-mail \textbf{sergio.novak@unirv.edu.br} ao final da prova. Pode-se utilizar o Google Colab para auxílio na implementação SQL.

\vspace{0.3cm}

Considere o seguinte esquema de banco de dados relacional para as questões de \textbf{1 a 5}:
\begin{itemize}[label={}]
    \item \textbf{tbEmpregado} (\underline{cod\_emp}, nome\_emp, rua\_emp, cidade\_emp, fone\_emp, cod\_gerente)
    \item \textbf{tbTrabalha} (\underline{cod\_emp}, \underline{cod\_companhia}, salário)
    \item \textbf{tbCompanhia} (\underline{cod\_companhia}, nome\_companhia, cidade\_companhia)
    \item \textbf{tbGerente} (\underline{cod\_gerente}, nome\_gerente)
\end{itemize}

\vspace{0.3cm}

\begin{enumerate}
    \item Escreva os comandos de definição de dados (\textbf{DDL}) para a criação das tabelas \textbf{tbGerente}, \textbf{tbCompanhia} e \textbf{tbEmpregado}.\\
    \textit{(Opcional: Resposta no PDF enviado por e-mail)}\\
    \resposta{2cm}

    \item Escreva o comando \textbf{DDL} para a criação da tabela \textbf{tbTrabalha}.\\
    \textit{(Opcional: Resposta no PDF enviado por e-mail)}\\
    \resposta{2cm}

    \item Implemente comandos de inserção (\textbf{INSERT}) para as tabelas.\\
    \textit{(Opcional: Resposta no PDF enviado por e-mail)}\\
    \resposta{2cm}

    \item Escreva o comando \textbf{DDL} (\textit{ALTER TABLE}) para adicionar a coluna \texttt{data\_nascimento}.\\
    \textit{(Opcional: Resposta no PDF enviado por e-mail)}\\
    \resposta{2cm}

    \item Escreva as consultas SQL:
    \begin{enumerate}[label=\alph*)]
        \item Liste o nome e a cidade de residência dos empregados da `Petrobras' com salário > 5000.
        \item Utilizando \textbf{LEFT JOIN}, liste todas as companhias e seus empregados.
    \end{enumerate}
    \textit{(Opcional: Resposta no PDF enviado por e-mail)}\\
    \resposta{2cm}
\end{enumerate}

\vspace{0.5cm} 


\end{document} 
